\documentclass[11pt,a4paper]{article}
\usepackage[margin=1in]{geometry}
\usepackage{hyperref}
\usepackage{enumitem}
\usepackage{graphicx}
\usepackage{xcolor}

% -----------------------------------------------------
% bvraghav-tiet-ucs503p-article.sty
% -----------------------------------------------------

% Variable: Lab Instructor
\makeatletter \newcommand{\BvrSetLabInstructor}[1]{%
  \gdef\Bvr@LabInstructor{#1}}
% default
\newcommand{\Bvr@LabInstructor}{Dr. Raghav %
  B. Venkataramaiyer}
% public accessor
\newcommand{\BvrLabInstructorValue}{\Bvr@LabInstructor}
\makeatother

% Add subtitle
\usepackage{titling}
% -- Set title bold
\pretitle{\begin{center}\bfseries\fontsize{18bp}{18bp}\selectfont}
\makeatletter
\newcommand{\BvrSubtitle}[1]{%
  \posttitle{%
    \par\end{center}
    \begin{center}\large#1\end{center}
    \vskip 0.5em}%
}
\makeatother

% Hints in the section title(s)
\newcommand{\BvrHint}[1]{%
  {\normalfont\normalsize\itshape\hfill #1}}

% -----------------------------------------------------

\title{SkinSense}

\BvrSetLabInstructor{Dr. Raghav B. Venkataramaiyer}

\BvrSubtitle{%
  Skin Analysis \& Skincare Recommendation Platform \\
  Project Proposal for UCS503P  \\
  Submitted to: \BvrLabInstructorValue \\ \bfseries
  Thapar Institute of Engineering and Technology%
}

\author{
  Muskan Kohli, CSED%
  \thanks{Roll No: \texttt{1024170453}, %
    Email: \texttt{mkohli1\_be24\@thapar.edu}}
  \and
  Himantveer Kaur, CSED%
  \thanks{Roll No: \texttt{1024170271}, %
    Email: \texttt{hkaur10\_be24\@thapar.edu}}
  \and
  Ansh Bindal, CSED%
  \thanks{Roll No: \texttt{1024170273}, %
    Email: \texttt{abindal\_be24\@thapar.edu}}
  \and
  Arnav Chandla, CSED%
  \thanks{Roll No: \texttt{1024170259}, %
    Email: \texttt{achandla\_be24\@thapar.edu}}
}

\date{\today}

\begin{document}
\maketitle

\section{Higher-order goal%
  \BvrHint{\ldots secondary framing}}
This project aims to make personalized skincare guidance
more accessible, reducing reliance on generic
marketing-driven advice or costly dermatologist
consultations for basic skin assessment. While the
broader goal is accessible self-care, the immediate
engineering objective is to translate a photo-based skin
assessment into a measurable, deployable software
solution that recommends suitable products and routines.

\section{Time-to-value %
\BvrHint{\ldots primary heuristic}}
\begin{itemize}[leftmargin=0.7cm]
    \item We will deliver meaningful increments quickly by getting the core pipeline (photo upload $\rightarrow$ classification $\rightarrow$ recommendation) working end-to-end early, so integration problems surface while there is still time to address them.
    \item We will prioritize fast feedback over perfection by using iterative development with CI-backed automation from the start, so every change is build-tested and deployable.
    \item Success is defined by what can be delivered and measured in the first iteration, then improved based on further testing.
\end{itemize}

\section{Problem Statement}
Most people do not have easy access to a dermatologist
for basic skin assessment, and end up relying on
subjective self-assessment or generic, marketing-driven
product recommendations that do not account for their
actual skin type, acne severity, or ingredient
allergies. Common pain points include:
\begin{itemize}[leftmargin=0.7cm]
    \item \textbf{Subjective self-assessment:} people often misjudge their own skin type (oily/dry/normal) or the severity of their acne, leading to mismatched product choices.
    \item \textbf{Generic recommendations:} most skincare advice available online is not personalized, and does not account for an individual's specific skin condition.
    \item \textbf{Ingredient conflicts and allergies:} users have no easy way to check whether a recommended product contains an ingredient they should avoid.
    \item \textbf{No structured routine:} even when good products are chosen, users lack guidance on how to sequence them into a coherent morning/evening routine.
\end{itemize}

\textbf{Why this matters now (contextual relevance):} personalized skincare guidance is currently gated behind either paid consultations or trial-and-error, and a system that gives structured, explainable recommendations from a single photo can meaningfully lower that barrier.

\section{Proposed Solution %
\BvrHint{\ldots Software Proposition}}
\subsection{Overview}
Build a \textbf{web-based} skin analysis and skincare
recommendation platform with the following components:
\begin{itemize}[leftmargin=0.7cm]
    \item \textbf{Photo upload and analysis:} users upload a face photo for skin type and acne severity assessment.
    \item \textbf{Skin type classification:} categorizes skin as oily, dry, or normal.
    \item \textbf{Acne severity classification:} a coarse-to-fine grading approach on the four-level Hayashi scale (mild/moderate/severe/very severe).
    \item \textbf{Rule-based recommendation engine:} matches skin type and acne severity against a curated product/ingredient database, with an allergy/ingredient exclusion filter and plain-language reasons for each recommendation.
    \item \textbf{Routine builder:} organizes recommended products into a morning/evening (AM/PM) routine.
    \item \textbf{Skin tone/undertone detection:} a lightweight, low-training-cost addition to refine recommendations (e.g., sunscreen suitability).
\end{itemize}

\subsection{Core workflow}
\begin{enumerate}[leftmargin=0.7cm]
    \item User uploads a face photo through the web interface.
    \item Backend runs the skin type and acne severity classifiers and returns their outputs.
    \item The recommendation engine matches these outputs, plus any allergy filters set by the user, against the product/ingredient database.
    \item The system returns a ranked product list with reasons, organized into an AM/PM routine.
\end{enumerate}

\subsection{Operational constraints for an educational setting}
\begin{itemize}[leftmargin=0.7cm]
    \item Keep the solution usable on low-to-mid range devices via responsive web design.
    \item Rely on established transfer-learning techniques, a published, benchmarked architecture for acne grading (Section 5.2), and deterministic rule-based logic elsewhere, avoiding unvalidated or ad hoc algorithms.
    \item Design for deployability using free-tier web hosting and managed databases.
\end{itemize}

\section{Solution Approach %
\BvrHint{\ldots Engineering focus}}
This is an engineering course project. The emphasis
throughout is on correct, testable system design and
integration, not on novel machine learning research.

\subsection{Recommendation logic %
\BvrHint{\ldots non-research}}
The recommendation engine is rule-based, not machine
learned:
\begin{itemize}[leftmargin=0.7cm]
    \item Match skin type, acne severity, skin tone, and allergy exclusions against tagged fields in the product/ingredient database.
    \item Apply deterministic matching rules sourced from the ingredient database's documented suitability tags, not invented ad hoc.
    \item Always show the reason a product was recommended or excluded; never silently drop a result.
\end{itemize}
This keeps the decision core deterministic and fully
unit-testable, which is deliberate: we have no reliable
ground-truth label for ``which product actually helped a
given user,'' so a learned recommender would not be
justifiable at this stage.

\subsection{Acne severity grading approach %
\BvrHint{\ldots engineering design}}
Since ACNE04 has far more mild/moderate images than
severe/very severe ones, a single 4-way classifier may
struggle on the rarer classes. As a first attempt, we
plan to try a two-step approach instead: a coarse
low-vs-high severity split, followed by a finer
classifier within each group. This is treated as an exploratory
attempt rather than a guaranteed outcome; if it does not
perform well in practice, a single classifier across all
four levels serves as the fallback.

\subsection{Web architecture %
\BvrHint{\ldots web-based solution}}
A typical 3-tier web architecture:
\begin{itemize}[leftmargin=0.7cm]
    \item \textbf{Frontend:} React, with a responsive upload page, results page, and a dark mode toggle.
    \item \textbf{Backend API:} FastAPI, handling image upload, model inference calls, and recommendation-engine queries.
    \item \textbf{Database:} PostgreSQL, storing the product catalog, ingredient data, and allergy/suitability tags, modeled with an explicit ER diagram.
\end{itemize}

\subsection{CI/CD and fast delivery enabler}
GitHub Actions will be used to:
\begin{itemize}[leftmargin=0.7cm]
    \item Automatically build and run backend unit tests and linting on every change.
    \item Deploy the application to a staging environment (Render/Railway) on successful builds.
    \item Enable safe, frequent releases of incremental features.
\end{itemize}

\section{Evaluation Criterion %
\BvrHint{\ldots measurable and attributable}}
We define success using measurable metrics tied to the
core classification and recommendation pipeline.

\subsection{Primary evaluation metrics}
\begin{itemize}[leftmargin=0.7cm]
    \item \textbf{Skin type classifier:} target $\geq$ 80\% accuracy on a held-out test set.
    \item \textbf{Acne severity classifier:} target $\geq$ 80\% accuracy, with a minimum acceptable floor of $\geq$ 70\%, given the added difficulty of this task. Macro-F1 is reported alongside accuracy, since acne severity classes are naturally imbalanced (mild/moderate cases far outnumber severe/very severe), and accuracy alone can mask poor performance on the rarer classes.
    \item \textbf{Attribution:} both metrics are computed against the held-out test split of the respective training datasets, using a fixed, documented split before any tuning begins.
\end{itemize}

\subsection{Secondary evaluation metrics}
\begin{itemize}[leftmargin=0.7cm]
    \item \textbf{Recommendation engine test coverage:} target 100\% coverage of rule branches, since the engine is deterministic and fully testable by design.
    \item \textbf{Inference time per scan:} measured and reported, to support a green-computing discussion of the lightweight backbone choice.
    \item \textbf{Allergy filter correctness:} verified via test cases that a flagged ingredient never appears in a returned recommendation.
    \item \textbf{User clarity:} pilot testers rate the clarity of the recommendation explanations on a 1--5 scale.
\end{itemize}

\subsection{Pilot validation plan %
\BvrHint{\ldots high-level}}
\begin{itemize}[leftmargin=0.7cm]
    \item Once the system is functional end-to-end, recruit a small group of classmates/friends to test the upload-to-routine flow. No testers are lined up as of this proposal; this is a stated target for the testing phase.
    \item Collect a 1--5 clarity rating plus open feedback on the recommendations and routine output.
\end{itemize}

\section{Scalability %
\BvrHint{\ldots with foresight}}
The proposition should scale in both theoretical and
practical senses.

\begin{enumerate}[leftmargin=0.7cm]
    \item \textbf{Scalable (theoretically):} the system should support growth in users and concurrent requests by:
    \begin{itemize}[leftmargin=0.7cm]
        \item decoupling frontend and backend components,
        \item indexing common product/ingredient queries in PostgreSQL,
        \item using a stateless API design.
    \end{itemize}

    \item \textbf{Operationally deployable (practically):} the system should run on common web hosting:
    \begin{itemize}[leftmargin=0.7cm]
        \item a managed PostgreSQL database,
        \item platform-hosted deployment (Render/Railway),
        \item a GitHub Actions CI/CD pipeline for staging.
    \end{itemize}
\end{enumerate}

\section{Engine availability heuristic}
A mature set of ``engines'' (tooling/services) is
available to implement this as a web-based project:
\begin{itemize}[leftmargin=0.7cm]
    \item Pretrained CNN backbones for transfer learning, via PyTorch: MobileNetV2 for skin type, ResNet-50 for the coarse-to-fine acne severity grader.
    \item Free GPU training environments (Google Colab / Kaggle Notebooks).
    \item Web framework for REST APIs (FastAPI) with auto-generated documentation.
    \item Managed relational database (PostgreSQL).
    \item Test framework for unit/integration tests (\texttt{pytest}).
    \item CI runner and staging deployment target (GitHub Actions, Render/Railway).
\end{itemize}
These components are already mature and well documented,
which lets us direct project effort toward model
integration, rule design, correctness, and measurement.

\section{Project scope and deliverables %
\BvrHint{\ldots iteration-friendly}}
\subsection{Initial deliverable %
\BvrHint{\ldots first iteration}}
\begin{itemize}[leftmargin=0.7cm]
    \item Photo upload interface
    \item Skin type classifier (oily/dry/normal) integrated behind a live API endpoint
    \item Acne severity classifier (4-level Hayashi grading; coarse-to-fine split attempted first, single classifier as fallback) integrated behind a live API endpoint
    \item Rule-based recommendation engine matching skin type and acne severity against the product database
    \item Basic frontend: upload page and results page
    \item CI: build + backend unit tests + linting
    \item CD to staging with a single pipeline job
\end{itemize}

\subsection{Subsequent deliverables}
\begin{itemize}[leftmargin=0.7cm]
    \item AM/PM routine builder
    \item Allergy/ingredient exclusion filter
    \item Skin tone/undertone detector
    \item ``Why recommended'' explanation tags for each product
    \item Lesion count estimation, decoupled from the severity classifier
\end{itemize}

\subsection{Stretch goals %
\BvrHint{\ldots Tier 3, attempted only if time remains}}
\begin{itemize}[leftmargin=0.7cm]
    \item A simple task queue (e.g., Celery) so a photo upload creates a job a worker processes asynchronously, with the frontend polling a status endpoint -- a basic, working version of event-driven design, not a production-grade one.
    \item Flagging low-confidence scans using Monte Carlo Dropout, so users can be asked to retake an unclear photo.
    \item Checking accuracy across different skin tones (Fitzpatrick17k), and trying a simple fix like balanced sampling if a real gap shows up.
    \item Trying the skin-type model in-browser (e.g., TensorFlow.js) and comparing its speed to the server-based version.
\end{itemize}
A multi-class lesion detector is noted as a possible
longer-term extension: fine-tune a pretrained
object-detection model (e.g., YOLO, pretrained on COCO)
on AcneSCU with moles, freckles, and melasma as
categories separate from acne, so these are not
misclassified as acne. This is not a committed
deliverable, and is contingent on data availability.

\section{Risks and mitigations}
\begin{itemize}[leftmargin=0.7cm]
    \item \textbf{Acne severity accuracy:} expected to score lower than skin type; the coarse-to-fine attempt (Section 5.2) is a first try at improving this, with a plain classifier as fallback, and a safety-floor target of $\geq$ 70\% either way.
    \item \textbf{Moles/freckles misread as acne:} mitigated by including non-acne examples in validation to measure the error, with a fine-tuned detector (Section 9.3) as a longer-term fix.
    \item \textbf{Photo lighting/angle variability:} plan basic preprocessing (cropping, lighting normalization) before classification.
    \item \textbf{ACNE04 access:} request from the original authors; a lower-quality but accessible Kaggle dataset is the fallback.
    \item \textbf{Evaluation measurement ambiguity:} define accuracy/F1 targets and a fixed test split before development begins.
    \item \textbf{Operational issues in staging:} deploy early; rely on CI/CD for repeatable deployments.
\end{itemize}

\section{Summary}
This project proposes a deployable web platform that
gives users a personalized, explainable skincare
routine from a single uploaded photo. It meets the
course criteria by emphasizing time-to-value through
iterative delivery, using CI/CD to support frequent
safe releases, and defining measurable evaluation
criteria for both the classification models and the
recommendation engine. The recommendation logic itself
is kept deterministic and rule-based by design, keeping
the project's focus on correct, testable engineering.

\end{document}
